\documentclass[11pt]{article}
\usepackage[T1]{fontenc}
\usepackage[utf8]{inputenc}
\usepackage{lmodern,amsmath,amssymb,amsthm,mathtools}
\usepackage[margin=25mm]{geometry}
\usepackage{microtype,longtable,array,booktabs,enumitem}
\usepackage[colorlinks=true,linkcolor=blue,urlcolor=blue]{hyperref}
\hypersetup{pdftitle={EMK Constructive Cut Foundation},pdfauthor={Monty Dabas research programme}}
\newtheorem{theorem}{Theorem}[section]
\newtheorem{proposition}[theorem]{Proposition}
\newtheorem{definition}[theorem]{Definition}
\newtheorem{corollary}[theorem]{Corollary}
\newcommand{\Hist}{\mathcal H}
\newcommand{\Qcut}{\mathbb Q_\Sigma}
\newcommand{\Rcut}{\mathbb R_\Sigma}
\newcommand{\Kcut}{\mathbb K_\Sigma}
\newcommand{\Id}{\mathrm I}
\newcommand{\End}{\operatorname{End}}
\newcommand{\Span}{\operatorname{span}}
\newcommand{\cut}{\kappa}
\newcommand{\selfcut}{\chi}
\setlength{\emergencystretch}{3em}
\setlist{nosep}
\title{EMK Constructive Cut Foundation\\
\large Typed records, native operators, iota, recognition geometry and completion}
\author{Research programme: Monty Dabas\\\small Corrective mathematical edition}
\date{1 October 2026}
\begin{document}
\maketitle
\begin{abstract}
We construct an explicit foundation from finite typed records of cut and
self-cut. Cut counts supply signed refinement scalars. The two recorded roles
supply parity and exchange maps whose composition has a derived four-cycle and
square minus identity; no complex field or cyclic phase alphabet is an input.
The complex-like scalar sector is subsequently recovered inside the larger
operator algebra. Finite recognition depth supplies idempotent apertures,
a prefix topology, a native ultrametric and a complete history space. Full
history remains distinct from operator or scalar shadow. Exact comparison-loop
residues characterize path-independent increments.

This is a \emph{corrected constructive theory}, not a proof of the original
manuscript's unrestricted three-prose-axiom claim. The typed recursion and
completion operations are stated explicitly. In particular, closure of the
constructed role operators is not literal closure of arbitrary primitive acts.
Every theorem in the new core has a written proof or an identified native
analytic dependency. Exact finite certificates and source hashes accompany it;
no proof-assistant verification or physical identification is claimed.
\end{abstract}
\bigskip

\section{The source, and what has been corrected}
\label{sec:source}
The input manuscript \texttt{emk\_topology.tex} has SHA-256
\begin{quote}\ttfamily\small
55f449e7eb3a2bde9e77193f595a6a5eeff3923da2bea076d00f0e9907b8cebd.
\end{quote}
It introduces a ground, a cut from the ground, and a self-cut on the trace of
a cut. Later it applies both acts recursively to all appearances, identifies
their product with an involution, and derives further structures from that
identification. Those are different mathematical assertions.

The present source is the \emph{free typed realization of the stated recursion}.
There are sorts Ground, History and Trace. A first-cut constructor gives a
seed $p_0$ from the ground. There is a history-extension constructor $k$;
recording a history gives its trace; a self-cut applied to that trace gives
the constructor $c$. Thus subsequent histories are finite words in the two
event tags $k,c$, based at $p_0$.

The history with no \emph{additional} event is denoted $\epsilon$.
It is the post-cut seed $p_0$, not the pre-cut ground and not absence of an
object. Zero additional cut count therefore does not identify the ground
with a scalar zero. The map from histories to their count is explicitly a
summary, not history equality.

The extension and trace constructors make the manuscript's recursive domain
use precise. They are a stated formalization, not a theorem that its original
prose uniquely forces this formalization. The proof language uses logical
equality, finite records, functions, subsets and induction. These are declared
metamathematical infrastructure. No ordinary complex numbers, spacetime,
metric, noise distribution or physical clock is an initial object.

\begin{theorem}[C1: Free record and evaluation theorem]
\label{thm:free}
Finite based records form a free monoid under concatenation. Every assignment
of the two tags to composable endomaps has a unique evaluation preserving
concatenation and the empty record. For general typed arrows, the analogous
construction has separate local identity paths and composition only at
matching endpoints.
\end{theorem}
\begin{proof}
The empty record must evaluate to the appropriate identity. Removing the last
tag of a nonempty record leaves a shorter record; its value and the value of
the final tag determine the evaluation. Induction gives existence and
uniqueness. Concatenation is associative because either parenthesization gives
the same ordered list. Typed endpoint matching proves the same statement
object by object.
\end{proof}

\begin{proposition}[C2: What primitive naming does not prove]
\label{thm:gap}
The free record $ckck$ is not the empty record. The equation
$(\selfcut\cut)^2=\Id$ is not a consequence of specifying two composable acts.
\end{proposition}
\begin{proof}
The two records have different lengths. Alternatively, take three appearances
and maps $\cut=(1,2,0)$, $\selfcut=(0,1,2)$ in table notation. The composite
squared sends $0$ to $2$, not $0$. A separate initial cut can still create
the seed $0$. Thus even the totalized algebraic reading permits failure of
the claimed equation.
\end{proof}
The result below derives closure on a \emph{constructed role carrier}. It does
not erase this counterexample or silently reinterpret a literal primitive
identity. This distinction is part of the certificate.

\section{Counts, signs and refinement before scalar coordinates}
\label{sec:counts}
Let a cut count be the length of a finite string of identical counting marks.
The empty counting string has count zero and appending a mark is successor.
Disjoint concatenation defines addition; replacing each mark by a copy of a
second counting string defines multiplication.

\begin{theorem}[C3: Signed and refined cut counts]
\label{thm:counts}
Signed pairs $(p,n)$ modulo
\[
(p,n)\sim(p',n')\quad\Longleftrightarrow\quad p+n'=p'+n
\]
form a commutative ordered ring of signed cut counts. Refinement packets
$[a:b]_\Sigma$, with signed $a$ and positive count $b$, modulo
$ab'=a'b$, form an ordered field $\Qcut$ under
\[
[a:b]+[u:v]=[av+ub:bv],\qquad [a:b][u:v]=[au:bv].
\]
\end{theorem}
\begin{proof}
Concatenation and rectangular replication give associative, commutative count
addition and multiplication, with distributivity, by matching finite positions.
Count cancellation follows by removing equal initial blocks, and induction
compares any two counts. Adding signed pairs componentwise and setting
$(p,n)(u,v)=(pu+nv,pv+nu)$ respects the displayed relation by expansion.
The swapped pair is the additive inverse. The sign of the difference is
well-defined by count comparison. A product of two positive nonzero differences
is positive, so there are no zero divisors.

Cross-multiplication shows the refinement relation is reflexive and symmetric;
transitivity uses cancellation of its positive common denominator. Expanding
the two operation formulas shows independence of representatives. The field
laws follow by putting finitely many operands over a common positive
denominator and using the signed ring laws. For nonzero $[a:b]$, interchange
numerator and denominator and move any denominator sign to the numerator;
the resulting product is one. Cross-multiplied signs give the field order.
\end{proof}
Opposite signs cancel in the \emph{coefficient ledger}. They do not delete the
underlying cut-event record or physically undo an irreversible cut. This is
the native count/refinement content credited to MP NR-7 in the source ledger.

\section{A four-cycle derived from the two cut roles}
\label{sec:roles}
Form the free signed ledger on the two distinct recorded roles
$e_k,e_c$. After refinement, denote it
$V=\Qcut e_k\oplus\Qcut e_c$. This is an additive construction on tags;
its elements are not asserted to be all UGD states.

Role exchange and role parity are the following explicitly constructed maps:
\[
K e_k=e_c,\quad K e_c=e_k,\qquad
H e_k=e_k,\quad H e_c=-e_c.
\]
Exchange is induced by the swap of the two tags. Parity takes the additive
inverse of the self-cut-role count, while leaving the cut-role count fixed.
Their extension is unique by additivity. The order of the named roles fixes
the displayed orientation convention; reversing it reverses the quarter-turn.

\begin{theorem}[C4: Native role-operator closure]
\label{thm:roles}
Define $R=KH$. Then, on every signed/refined role ledger,
\[
H^2=K^2=\Id,\qquad R^2=-\Id,\qquad R^4=\Id,
\qquad KR=-RK.
\]
The orbit of $e_k$ under $R$ is precisely
$e_k,e_c,-e_k,-e_c,e_k$. The group generated by $H,K$ consists of the
eight signed permutations $\pm\Id,\pm H,\pm K,\pm R$.
\end{theorem}
\begin{proof}
Directly on the two free generators,
$Re_k=e_c$ and $Re_c=-e_k$. Applying twice negates each generator, and
applying four times fixes each. Exchange and parity each square to identity.
Also $KRe_k=e_k=-RKe_k$ and $KRe_c=-e_c=-RKe_c$, proving anticommutation.
Additivity extends each identity to every ledger.

All products reduce to the eight listed signed permutations using these
relations. Their images of the two generators are pairwise distinct; each
listed permutation is generated by $H,K$. Thus the list is exact.
\end{proof}
Here the number four is an \emph{output}, not a supplied phase-alphabet size.
The constructed maps are definitions on the free role completion; the
closure equations are their proved consequences. No theorem has identified
these maps with every possible physical realization of primitive acts.

The free evaluation of Theorem~\ref{thm:free} permits the assignment
$k\mapsto H$, $c\mapsto K$. Its image of the composite $ck$ is $R$.
Under this particular interpretation $(ck)^2$ acts as $-\Id$, and
$(ck)^4$ acts as $\Id$. Four quarter-turns therefore mean eight alternating
primitive-tag occurrences here. They are not the upload's asserted four
primitive-tag identity. This correction prevents an off-by-two cycle claim.

\section{The larger operator algebra and its scalar sector}
\label{sec:scalars}
\begin{theorem}[C5: Algebra and recovered scalar sector]
\label{thm:scalars}
The algebra generated by the role operators is all $\End_{\Qcut}(V)$, with
basis $\Id,K,R,RK$. The operators commuting with $R$ are precisely
\[
\Kcut^0=\{a\Id+bR:a,b\in\Qcut\}.
\]
This commutative sector is a field, with derived element $\iota=R$,
\[
\iota^2=-1,\quad
(a+b\iota)(c+d\iota)=(ac-bd)+(ad+bc)\iota,
\quad
(a+b\iota)^{-1}=\frac{a-b\iota}{a^2+b^2}
\]
for nonzero $a+b\iota$.
\end{theorem}
\begin{proof}
Read the two generator images after constructing the operators. Their output
coordinate arrays are
\[
\Id=\begin{pmatrix}1&0\\0&1\end{pmatrix},\quad
K=\begin{pmatrix}0&1\\1&0\end{pmatrix},\quad
R=\begin{pmatrix}0&-1\\1&0\end{pmatrix},\quad
RK=\begin{pmatrix}-1&0\\0&1\end{pmatrix}.
\]
They span all four independent array positions because division by two exists
in $\Qcut$. Comparing the two images in $AR=RA$ forces an arbitrary array
$A$ to have equal diagonal entries and opposite off-diagonal entries, hence
$A=a\Id+bR$. Its product law follows from $R^2=-\Id$. It is closed and
commutative. The positive ordered-cut expression $a^2+b^2$ vanishes only
when both coefficients vanish. Multiplying by the displayed inverse gives
identity. Associativity and distributivity come from composition and addition
of the constructed maps.
\end{proof}
These arrays are outputs of the free-generator calculation. They were not
assumed as primitive geometry. Further, $R$ is central in $\Kcut^0$ but not
in the full algebra: it anticommutes with $K$. The scalar sector therefore
cannot replace the noncommuting algebra or the full UGD state.

\section{Pairing, dagger, seam and operator curvature}
\label{sec:pairing}
Coefficient matching of the two recorded roles defines
\[
B(ae_k+be_c,ue_k+ve_c)=au+bv,\qquad
\omega(ae_k+be_c,ue_k+ve_c)=av-bu.
\]
\begin{theorem}[C6: Native pairing and compatible orientation]
\label{thm:pairing}
$B$ is bilinear, symmetric and positive definite. It is preserved by $H,K,R$.
The oriented form satisfies $\omega(x,y)=B(Rx,y)$. Within the scalar sector,
orientation reversal gives
\[
(a+b\iota)^\dagger=a-b\iota,\quad
z^\dagger z=a^2+b^2,\quad
N(zw)=N(z)N(w).
\]
\end{theorem}
\begin{proof}
Bilinearity and symmetry follow by distributivity. Its diagonal is the sum
of the two signed-cut squares, positive for a nonzero ledger. Sign changes
and tag exchange preserve matched products, hence so does their composite.
Substitution of $Rx=-be_k+ae_c$ gives the oriented-form identity. Reversing
the quarter-turn changes $R$ to $-R$; expansion proves the dagger laws and
the norm-square identity. Expanding the product norm gives
$(a^2+b^2)(c^2+d^2)=(ac-bd)^2+(ad+bc)^2$.
\end{proof}
The pairing is derived from this marked-role counting construction. It is
not declared to be a universal physical spacetime metric, nor automatically
the pairing on every future UGD or morphic module.

\begin{theorem}[C7: Seam aperture and genuine mixed curvature]
\label{thm:seam}
The fixed seam of $K$ is $\Qcut(e_k+e_c)$, including both signs. Put
$P=(\Id+K)/2$ and $Q=(\Id-K)/2$. Then
\[
P^2=P,\quad Q^2=Q,\quad PQ=QP=0,\quad P+Q=\Id.
\]
The native mixed order defect is $[R,K]=2RK\ne0$, whereas $[K,K]=0$.
For any two endomaps $A,B$ of $V$,
\[
PABP-(PAP)(PBP)=PAQBP.
\]
\end{theorem}
\begin{proof}
Exchange fixes exactly equal role coefficients. Expand the projector
expressions using $K^2=\Id$. Anticommutation gives the mixed defect; its
value on $e_k$ is $-2e_k$, which is nonzero. The final identity follows by
inserting $\Id=P+Q$ between $A$ and $B$.
\end{proof}
The last formula credits the existing native cut-corner theorem. It identifies
the returning omitted component algebraically, without assigning a probability
law or physical noise mechanism. That later work remains separate.

\section{Cut-Cauchy scalar completion and native analytic dependencies}
\label{sec:analytic}
Define $D(a+b\iota)=|a|+|b|$ in the ordered cut field. Direct expansion
proves the triangle inequality and $D(zw)\le D(z)D(w)$.

\begin{theorem}[C8: Complete native coefficient field]
\label{thm:completion}
Cut-Cauchy sequences in $\Kcut^0$, modulo differences tending to zero in
every positive rational-cut radius for $D$, form a complete field $\Kcut$.
Termwise addition, multiplication and dagger are well-defined, and the
constant-sequence map embeds $\Kcut^0$ densely. The radial component is
the ordered Archimedean completion $\Rcut$ of $\Qcut$. Each nonnegative
radial scalar has a unique nonnegative square root, and $\sqrt{N(z)}$
is a multiplicative norm defining the same completion.
\end{theorem}
\begin{proof}
A Cauchy sequence is bounded: after its Cauchy threshold it stays within one
unit of a fixed term, and finitely many earlier terms have a common count
bound. The two inequalities above show that sums and products of Cauchy
sequences are Cauchy and that changing either by a null sequence changes
the result by a null sequence. Dagger preserves $D$.

A nonzero Cauchy class has some positive rational $\epsilon$ for which
arbitrarily late terms have $D\ge\epsilon$. Choose a Cauchy tail whose
pairwise differences have $D<\epsilon/2$. Comparison with a later large
term shows every term of that tail has $D\ge\epsilon/2$. Write this lower
bound as $\delta$. Since $a^2+b^2\ge(|a|+|b|)^2/2$, the inverse formula
gives $D(z^{-1})\le2/\delta$. Consequently
\[
D(z_n^{-1}-z_m^{-1})\le(4/\delta^2)D(z_n-z_m).
\]
Tail inverses are Cauchy and give the inverse class; finitely many initial
choices do not affect it.

Every class is approximated by sufficiently late individual terms of a
representative, so constant cut scalars are dense. For completeness,
approximate the $j$th member of a Cauchy sequence of
classes by a rational-cut element within $2^{-j}$. Enumerate rational-cut
pairs by increasing total numerator/denominator count, breaking ties by
the ordered role tags. Density supplies an eligible element; select the
first eligible element of this fixed enumeration. Here radius comparison
between classes means eventual comparison of representative differences,
with strict rational slack, so it is independent of representatives. No
simultaneous choice of representatives is required. The approximating
elements are Cauchy by the triangle inequality and their class is the limit.
Constant sequences are injective because a
nonzero exact cut scalar has a positive rational separation from zero.

Radial sequences remain radial under these operations. Their nonnegative
cone is the closure of the nonnegative rational cuts; comparison by rational
separation gives a proper total order. A Cauchy representative is bounded
by a finite cut count, proving the Archimedean property. Componentwise
construction identifies the full completed field with two radial components
and the already derived multiplication law, without taking ordinary
$\mathbb C$ as input.

For a nonnegative radial $s$, choose a count $M>s+1$. Bisect $[0,M]$,
keeping the half whose squared endpoints bracket $s$. The nested endpoints
are Cauchy because their separation is $M2^{-n}$. Their common limit $u$
satisfies $u^2=s$, since products respect limits. If $0\le u<v$, then
$v^2-u^2=(v-u)(v+u)>0$, proving uniqueness. Thus square root is constructed,
not imported. Finally
$(ac+bd)^2\le(a^2+b^2)(c^2+d^2)$ follows by subtracting to obtain
$(ad-bc)^2\ge0$. Expanding $N(z+w)$ and taking the nonnegative roots gives
the triangle inequality for $\sqrt N$. Positivity and multiplicativity
follow from Theorem~\ref{thm:pairing}. The inequalities
$N\le D^2\le2N$ prove agreement of their Cauchy notions.
\end{proof}
This is the constructive content of MP NR-7/NR-8; the original oriented-plane
input has here been supplied by Theorem~\ref{thm:roles}. Norm-square and $D$
neighborhoods agree because $N\le D^2\le2N$, so their completions coincide.

The following \emph{native analytic results are cited, not re-proved or
presented as a new formal certification here}: RKF F00/E (N1) constructs the
factorial-series exponential and Euler identities on this completed field;
F00G (N2) constructs the positive logarithm and proves its product law; RH F00J
(N3) constructs the fundamental circular period, and F00K (N4) proves scalar polar
decomposition and the branch logarithm. Their exact files, commits and
premises are in the source ledger. The coefficient and completion inputs
are supplied above. They do not select a physical angle trajectory or make
a scalar logarithm single-valued on full loop histories.

For clarity, the radial monotonicity input of F00G needs no additional
classical calculus result: for $t>0$ the native exponential series gives
$\operatorname{Exp}_\Sigma(t)>1$. The exponential group law gives positivity
also for negative $t$ by reciprocation. For $y>x$, it then gives
$\operatorname{Exp}_\Sigma(y)=\operatorname{Exp}_\Sigma(x)
\operatorname{Exp}_\Sigma(y-x)>\operatorname{Exp}_\Sigma(x)$.

\section{Recognition depth, topology and a native history metric}
\label{sec:history}
For a finite history $h$, read its two-tag word followed by an end marker
$\bot$, repeated after termination. The marker is a record of termination,
not the pre-distinction ground. Recognition at depth $n$ reads the first
$n$ symbols of this padded record. Two histories are fully recognized as
equal exactly when every finite readout agrees.

\begin{theorem}[C9: Recognition equality and prefix topology]
\label{thm:topology}
Full finite-depth recognition is equality of histories. Sets fixing a finite
readout form a clopen basis for a topology $\tau_{\rm rec}$. This is the
coarsest topology in which every finite recognition map into its finite
discrete record set is continuous.
\end{theorem}
\begin{proof}
Unequal words differ at a symbol or at the first termination of the shorter
word, so some finite readout separates them. A finite readout cylinder is
partitioned by the next-symbol cylinders. Intersections of compatible
cylinders are the longer cylinder; incompatible intersections are empty.
This proves the basis laws. The complement of a cylinder is a union of the
other cylinders at that depth, so it is clopen. Any topology making those
finite maps continuous contains every cylinder; their union/intersection
closure proves the claimed minimality.
\end{proof}
On finite histories this topology is discrete: reading through a finite
history's end marker isolates it. The same cylinders have a nontrivial
boundary topology on completed infinite histories below. The forward-invariant
topology of cut reachability is another construction and must not be
identified with this recognition topology.

\begin{theorem}[C10: Derived prefix ultrametric]
\label{thm:metric}
Let $\ell(h,g)$ be the length of the common prefix before the first
differing tag or end marker. Set
\[
d(h,g)=\begin{cases}0&h=g,\\2^{-\ell(h,g)}&h\ne g.\end{cases}
\]
Then $d$ is a metric with the stronger triangle inequality
\[
d(h,u)\le\max\{d(h,g),d(g,u)\}.
\]
It generates $\tau_{\rm rec}$. Prepending the same cut halves nonzero
distance; appending the same cut is non-expansive.
\end{theorem}
\begin{proof}
Symmetry and separation follow from first difference. If each of two pairs
agrees on its first $n$ symbols, so does the third pair. Hence
$\ell(h,u)\ge\min(\ell(h,g),\ell(g,u))$, which is the displayed inequality.
A ball $d(h,g)<2^{-n}$ fixes symbols through index $n$. Conversely every
finite cylinder contains a ball at each of its members. Thus the topologies
agree. Prepending adds one to the common-prefix length. Appending cannot
destroy a previously matching prefix, including when one original word ends
first; hence it cannot increase distance.
\end{proof}
The factor $1/2$ uses the native halving refinement of cut counts; it is an
explicit normalization, not a number forced uniquely by the tag alphabet.
More generally any rational $0<r<1$ in place of $1/2$ yields the same prefix
topology. The numerical calibration is therefore a construction convention,
not a selected physical length scale. This metric measures recorded
distinguishability; it is not claimed to be Riemannian curvature geometry.

\section{History completion, Eye apertures and the tower}
\label{sec:history-completion}
Adjoin infinite two-tag histories to the finite histories. A completed record
either has no end marker, or from its first end marker onward has only end
markers. Extend the prefix metric by the same first-difference rule.

\begin{theorem}[C11: Exact history completion]
\label{thm:history-completion}
The resulting space $\overline{\Hist}$ is complete, and finite histories are
dense. It is the metric completion of the finite source records.
\end{theorem}
\begin{proof}
In a Cauchy sequence, each finite prefix eventually stabilizes: take a radius
smaller than the corresponding dyadic separation. The stabilized prefixes
are compatible. If an end marker occurs, no later symbol can be a cut tag,
because every sufficiently late original record obeys this finite condition.
They therefore specify exactly one permitted finite or infinite history.
Prefix stabilization is precisely convergence to it. Every infinite history
is approximated within $2^{-n}$ by its finite length-$n$ truncation; finite
ones already belong to the dense source. The same description shows that
equivalent Cauchy sequences have the same limit profile, establishing the
completion identification directly.
\end{proof}

\begin{theorem}[C12: A correctly typed Eye hierarchy]
\label{thm:eye}
Let $E_n:\overline{\Hist}\to\overline{\Hist}$ truncate a history after at
most $n$ tags, retaining the result as an actual finite record. Then
\[
E_n^2=E_n,\qquad E_mE_n=E_{\min(m,n)},\qquad
d(E_nh,h)\le2^{-n}.
\]
Every history is recovered uniquely from its coherent truncations. Each
$E_n$ is non-expansive. There is no finite $n$ that recovers every history.
\end{theorem}
\begin{proof}
Truncating twice keeps the shorter limit, proving the two algebraic identities.
The truncated and original records agree in the first $n$ positions unless
already equal, proving the error bound. Identical coherent truncations fix
every symbol, hence the whole history; the bound proves convergence back to
that history. Truncation preserves any common prefix up to its cutoff and
otherwise makes the results equal, proving non-expansiveness. Histories
$k^n$ and $k^nc$ have equal $E_n$ and are distinct.
\end{proof}
These are endomaps with a common domain, so idempotence is genuinely typed.
For $m>n$, $E_m\ne E_n$ globally, as witnessed by a sufficiently long history.
A particular short history may stabilize. Thus the global hierarchy can be
strict without falsely requiring every individual state to change at every
level. It is a refinement tower, not powers of one involution.

\section{State return, history return and lifted phase}
\label{sec:memory}
\begin{theorem}[C13: No loss of cut history at operator return]
\label{thm:memory}
An evaluated operator word can be identity while its based event record is
nonempty. For the derived quarter-turn, four $R$ occurrences return every
role ledger but leave four operation records. For $n$ positive quarter-turn
occurrences there is a unique decomposition $n=4w+r$, $0\le r<4$, with
composition law
\[
(r,w)(s,v)=((r+s)\bmod4,\ w+v+\lfloor(r+s)/4\rfloor).
\]
\end{theorem}
\begin{proof}
The role return follows from $R^4=\Id$; a four-letter word is not the empty
word by Theorem~\ref{thm:free}. Repeatedly remove complete four-mark blocks
to obtain the unique remainder and count. Adding two counts gives the formula,
whose associativity is inherited from count addition. Forgetting $w$ sends
the counts zero and four to the same phase but erases their difference.
\end{proof}
The full word contains more than its loop count. UGD phase/scale/seam objects
must retain their active memory rather than be replaced by the scalar sector
of Theorem~\ref{thm:scalars}. MP NR-9 proves the appropriate scalar-shadow
quotient for its declared UGD law; it does not equate that quotient with the
full UGD state. No identification of all history words with a single seam
integer is made here.

Cut count gives an internal event index. Histories form a prefix order in
which extension increases that count. Different branches need not be ordered
as one physical time line. Neither event count nor a finite endpoint orbit
establishes a continuous physical clock or a helical spacetime manifold.

\section{Complete path-equality criterion from comparison residues}
\label{sec:paths}
Let a finite cut graph have oriented edges and a
signed/refined increment $a_e$ on each edge. Formal reversed traversal carries
$-a_e$. Reversal is a \emph{comparison of records}, not a physical reverse
operation that erases a cut. Choose a spanning forest by exploring finite
paths. Put a potential $F$ at zero on each root and integrate increments
along the unique tree paths.

\begin{theorem}[C14: Exactness and the full loop ledger]
\label{thm:paths}
For every edge $e:x\to y$ define $r_e=a_e-(F(y)-F(x))$.
Tree-edge residues vanish. The following are equivalent:
\begin{enumerate}[label=(\roman*)]
\item every edge increment is a potential difference;
\item the increment sum of every closed comparison walk is zero;
\item every non-tree residue $r_e$ is zero.
\end{enumerate}
There are $|E|-|V|+b$ independent fundamental comparison cycles, where $b$
is the number of connected components. Same endpoints alone do not force
equal increments.
\end{theorem}
\begin{proof}
Potential differences telescope on every walk, proving (i) implies (ii).
The tree path from a root to $x$, followed by $e$, then the reversed tree
path from $y$, has sum $r_e$, proving (ii) implies (iii). By the definition
of $F$, tree increments are already its differences, so (iii) gives (i).

Each non-tree edge yields one such cycle with coefficient one on that edge
and zero on other non-tree edges; hence these cycles are independent. Subtract
these cycles from any signed cycle to leave a cycle supported on a forest.
At a leaf its incident coefficient must be zero. Repeated leaf removal proves
the remainder zero. A forest has $|V|-b$ edges, giving the stated count.
\end{proof}
This supplies a complete finite criterion, including global periods; checking
only one local commutator is insufficient. A directed diamond has no directed
cycle but can have unequal sums on its two routes. A four-edge phase loop
with unit increments has lifted sum four despite endpoint return. In a later
native logarithm chart the analogous loop period must likewise be retained.
For the graph actually generated by the eight signed role operators and
left multiplication by $H,K$, there are eight vertices and sixteen directed
edges, hence nine independent comparison cycles. Its potential basis and
every chord defect are included in the certificate.

\section{What the completed core does and does not certify}
\label{sec:status}
The new core consists of Theorems/Propositions C1--C14 above. Each has a
stated construction and proof. Its mathematical inputs are the typed free
cut-record realization, the declared proof language and the explicit
completion operations. The analytic extensions in Section~\ref{sec:analytic}
are named existing native results with their own evidence levels.

The original claim that \emph{arbitrary} acts named first-cut and self-cut
force all later structures remains false or unproved as recorded in the
claim ledger. The new construction is not logically interchangeable with
that universal statement. In particular:
\begin{itemize}
\item the four-cycle is proved on the constructed signed role carrier;
\item the source's asserted raw $(\selfcut\cut)^2=\Id$ is not restored;
\item topology and metric are the constructed recognition geometry, not an
  automatic thermodynamic manifold;
\item Maxwell relations, Legendre potentials, Onsager response, physical
  event probabilities and an infinite helix are not certified from a graph
  cycle or operator notation;
\item the full UGD seam state is not equated with a complex scalar;
\item no RH, Yang--Mills continuum or physical-constant claim is added.
\end{itemize}
Removing an incorrect theorem from the active core is a documented correction,
not proof of the original theorem. The original file and all 48 dispositions
remain available alongside this edition. Thus \emph{the corrected constructive
core has a closed proof ledger}; \emph{the entire original manuscript is not
certified as written}. Physical and later extra-work connections are deferred.

\section{Certificate and source ledger}
\label{sec:certificate}
Run \texttt{certificate/verify.py} with the pinned RKF checkout. The local
checker exhausts the eight derived signed permutations, all 512 group triples
and all 32 pairing basis cases. These basis checks cover the full linear
identities, not only sampled coefficients. It also tests 63 finite histories,
all 250047 ultrametric triples, truncation bounds, 270 exact potential
assignments on four graph types, and negative controls. General infinite
statements have the written proofs above; the checker does not pretend to
enumerate all infinite histories or analytic points.

The unchanged canonical RKF engine supplies algebra normalization and replay.
The F00/E and F00GHI certificates are rerun against original source pins;
their analytical theorems are not thereby converted into proof-assistant
certificates. No Lean/Coq runtime is available in this environment, and no
formal-proof-assistant verification is claimed. Source hash, written proof,
finite algorithmic evidence and external review remain distinct.

Primary sources, with premises attached:
\begin{enumerate}
\item MP, commit \texttt{feedfca094aa557086a612a0736aaffc07ab1ae4}.\par
\path{proof_lab/native_operator_algebra/}: NR-7 signed cut/refinement
field; NR-8 cut-Cauchy completion; NR-9 UGD scalar-shadow quotient.
Native algebra/completion results. NR-7 originally \emph{defines} the
oriented action before proving its square; this edition constructs it from
role exchange and parity. Private source bodies are cited, not republished.
\item RKF, commit \texttt{3cc5a33b05c16d59c90994ddda69dedc0d392424},
\texttt{operator\_foundation/}: native typed paths, exact cut arithmetic,
normal-form compiler, faithful regular action and cut-corner residue.
Native presentation inputs; no primitive ordinary Hilbert space. Canonical
engine source is reused unchanged, not copied into this packet.
\item Same RKF commit, F00/E and F00G: native factorial-series Euler
construction and positive logarithm. Inputs: the ordered cut completion and
derived iota. Their finite source certificates are replayed; continuous
results remain supported by the source's written proofs.
\item RH, commit \texttt{c588ded973a395b5fce37c82f670616160979a2e}, F00J
and F00K: native period and scalar polar/logarithm results. Inputs: native
Euler functions and ordered completion. Written-source citation, no new
standalone formal or numerical certificate asserted here.
\item MP NR-13 normalized UGD group: algebraic gauge normalization is exact
for its declared potential cocycle. Its later ordinary Haar and $L^1/L^2$
extension invokes analytic group/measure machinery and Young-type bounds.
Those \emph{analytic imports are not dependencies of this core}.
\end{enumerate}
Exact URLs, Git blobs, SHA-256 hashes and evidence roles are recorded in
\path{certificate/SOURCE_PINS.json}. Hashing establishes source identity,
not theorem truth.

\appendix
\section{Disposition of every original claim}
The machine-readable ledger retains original environment titles and line
ranges. A replacement is a corrected statement, not retroactive approval of
the original wording. The following compact table covers all 48 original
theorem, lemma, proposition and corollary environments, including duplicates.
\small
\begin{longtable}{@{}>{\raggedright\arraybackslash}p{0.08\linewidth}>{\raggedright\arraybackslash}p{0.20\linewidth}>{\raggedright\arraybackslash}p{0.66\linewidth}@{}}
\toprule
Claim & Disposition & Corrected scope \\
\midrule
\endfirsthead
\toprule
Claim & Disposition & Corrected scope (continued) \\
\midrule
\endhead
\bottomrule
\endfoot
ET01 & Typed replacement & Free typed constructors make recursion lawful; the upload did not supply their recursive domains. \emph{C1.} \\[5pt]
ET02 & Recognition replacement & Equality of all finite cut-record observations is an explicit equivalence. Observer wording alone does not define it. \emph{C9.} \\[5pt]
ET03 & Recognition replacement & Equality of all finite cut-record observations is an explicit equivalence. Observer wording alone does not define it. \emph{C9.} \\[5pt]
ET04 & Recognition replacement & Equality of all finite cut-record observations is an explicit equivalence. Observer wording alone does not define it. \emph{C9.} \\[5pt]
ET05 & Recognition replacement & Equality of all finite cut-record observations is an explicit equivalence. Observer wording alone does not define it. \emph{C9.} \\[5pt]
ET06 & Typed Eye replacement & Truncations are endomaps of one completed history space and satisfy the proved aperture meet law. \emph{C12.} \\[5pt]
ET07 & Seam corrected & The actual exchange fixed set is the whole signed diagonal; its projector is derived. The original zero-cost and positive-ray statements are not certified. \emph{C7.} \\[5pt]
ET08 & Topology replaced & Recognition cylinders, their prefix metric and completion are constructed. They are not the upload's unsaturated image-stability topology. \emph{C9, C10, C11.} \\[5pt]
ET09 & Event index only & Cut count is an internal record index; a physical clock is not derived. \emph{C13.} \\[5pt]
ET10 & Role closure corrected & Exchange/parity produce a quarter-turn on the signed role carrier. Arbitrary raw chi/kappa closure and a four-point orbit of one involution do not follow. \emph{C2, C4.} \\[5pt]
ET11 & Tower corrected & Distinct truncation maps give a strict global hierarchy; powers of an involution do not. Individual finite histories can stabilize. \emph{C12.} \\[5pt]
ET12 & Completion constructed & Signed cut refinements and Cauchy completion construct scalars; arbitrary finite/infinite histories form a separate completed record space. Alternation alone is insufficient. \emph{C3, C8, C11.} \\[5pt]
ET13 & Typed Eye replacement & Truncations are endomaps of one completed history space and satisfy the proved aperture meet law. \emph{C12.} \\[5pt]
ET14 & Seam corrected & The actual exchange fixed set is the whole signed diagonal; its projector is derived. The original zero-cost and positive-ray statements are not certified. \emph{C7.} \\[5pt]
ET15 & Topology replaced & Recognition cylinders, their prefix metric and completion are constructed. They are not the upload's unsaturated image-stability topology. \emph{C9, C10, C11.} \\[5pt]
ET16 & Pairing and order split & A native role pairing is derived, but distinct role operators need not commute; the mixed defect is nonzero. \emph{C4, C6, C7.} \\[5pt]
ET17 & Full residue criterion & Endpoint equality does not erase periods. Potential exactness is equivalent to vanishing on every fundamental comparison cycle. \emph{C13, C14.} \\[5pt]
ET18 & Infinite sum withdrawn & The scalar/history completions and finite loop criterion are proved. They do not validate the unspecified infinite master-seam series. \emph{C8, C11, C14.} \\[5pt]
ET19 & Native analytic citation & Euler, radial logarithm, period and scalar polar/logarithm use pinned native theorems. The upload's arbitrary analytic-class equivalence is withdrawn. \emph{C4, C5, C8, N1, N2, N3, N4.} \\[5pt]
ET20 & Mixed defect corrected & The mixed operator commutator is nonzero; every self-commutator is zero. This does not identify it with Riemannian or connection curvature. \emph{C7, C14.} \\[5pt]
ET21 & Full residue criterion & Endpoint equality does not erase periods. Potential exactness is equivalent to vanishing on every fundamental comparison cycle. \emph{C13, C14.} \\[5pt]
ET22 & Native analytic citation & Euler, radial logarithm, period and scalar polar/logarithm use pinned native theorems. The upload's arbitrary analytic-class equivalence is withdrawn. \emph{C4, C5, C8, N1, N2, N3, N4.} \\[5pt]
ET23 & Role closure corrected & Exchange/parity produce a quarter-turn on the signed role carrier. Arbitrary raw chi/kappa closure and a four-point orbit of one involution do not follow. \emph{C2, C4.} \\[5pt]
ET24 & Full residue criterion & Endpoint equality does not erase periods. Potential exactness is equivalent to vanishing on every fundamental comparison cycle. \emph{C13, C14.} \\[5pt]
ET25 & Infinite sum withdrawn & The scalar/history completions and finite loop criterion are proved. They do not validate the unspecified infinite master-seam series. \emph{C8, C11, C14.} \\[5pt]
ET26 & Native analytic citation & Euler, radial logarithm, period and scalar polar/logarithm use pinned native theorems. The upload's arbitrary analytic-class equivalence is withdrawn. \emph{C4, C5, C8, N1, N2, N3, N4.} \\[5pt]
ET27 & Topology replaced & Recognition cylinders, their prefix metric and completion are constructed. They are not the upload's unsaturated image-stability topology. \emph{C9, C10, C11.} \\[5pt]
ET28 & Topology replaced & Recognition cylinders, their prefix metric and completion are constructed. They are not the upload's unsaturated image-stability topology. \emph{C9, C10, C11.} \\[5pt]
ET29 & Typed maps retained & Ordinary typed composition is valid. New recognition is full record equality, and truncations retain their common codomain; no unspecified quotient descent is assumed. \emph{C1, C9, C12.} \\[5pt]
ET30 & Typed maps retained & Ordinary typed composition is valid. New recognition is full record equality, and truncations retain their common codomain; no unspecified quotient descent is assumed. \emph{C1, C9, C12.} \\[5pt]
ET31 & Path and action separated & A free nonempty path may evaluate to identity on the role carrier while its event record remains distinct. \emph{C1, C4, C13.} \\[5pt]
ET32 & Cycle scope corrected & A specified directed cycle has only empty/full forward-invariant subsets. Its carrier is not its one-point orbit quotient, and that invariant topology differs from recognition topology. \emph{C4, C9.} \\[5pt]
ET33 & Cycle scope corrected & A specified directed cycle has only empty/full forward-invariant subsets. Its carrier is not its one-point orbit quotient, and that invariant topology differs from recognition topology. \emph{C4, C9.} \\[5pt]
ET34 & Physics claim withdrawn & Finite comparison potentials and their loop criterion are mathematical results; Maxwell, Legendre and thermodynamic identifications have not been derived. \emph{C14.} \\[5pt]
ET35 & Path and action separated & A free nonempty path may evaluate to identity on the role carrier while its event record remains distinct. \emph{C1, C4, C13.} \\[5pt]
ET36 & Cycle scope corrected & A specified directed cycle has only empty/full forward-invariant subsets. Its carrier is not its one-point orbit quotient, and that invariant topology differs from recognition topology. \emph{C4, C9.} \\[5pt]
ET37 & Physics claim withdrawn & Finite comparison potentials and their loop criterion are mathematical results; Maxwell, Legendre and thermodynamic identifications have not been derived. \emph{C14.} \\[5pt]
ET38 & Physics claim withdrawn & Finite comparison potentials and their loop criterion are mathematical results; Maxwell, Legendre and thermodynamic identifications have not been derived. \emph{C14.} \\[5pt]
ET39 & Physics claim withdrawn & Finite comparison potentials and their loop criterion are mathematical results; Maxwell, Legendre and thermodynamic identifications have not been derived. \emph{C14.} \\[5pt]
ET40 & Role closure corrected & Exchange/parity produce a quarter-turn on the signed role carrier. Arbitrary raw chi/kappa closure and a four-point orbit of one involution do not follow. \emph{C2, C4.} \\[5pt]
ET41 & Seam corrected & The actual exchange fixed set is the whole signed diagonal; its projector is derived. The original zero-cost and positive-ray statements are not certified. \emph{C7.} \\[5pt]
ET42 & Inverse scope corrected & The constructed role maps have proved inverses; arbitrary primitive acts are not declared invertible. \emph{C1, C4.} \\[5pt]
ET43 & Role closure corrected & Exchange/parity produce a quarter-turn on the signed role carrier. Arbitrary raw chi/kappa closure and a four-point orbit of one involution do not follow. \emph{C2, C4.} \\[5pt]
ET44 & Tower corrected & Distinct truncation maps give a strict global hierarchy; powers of an involution do not. Individual finite histories can stabilize. \emph{C12.} \\[5pt]
ET45 & Helix claim withdrawn & Infinite records, a refinement tower and retained winding are constructed. A physical helix or torus is not forced by an involution. \emph{C11, C12, C13.} \\[5pt]
ET46 & Mixed defect corrected & The mixed operator commutator is nonzero; every self-commutator is zero. This does not identify it with Riemannian or connection curvature. \emph{C7, C14.} \\[5pt]
ET47 & Topology replaced & Recognition cylinders, their prefix metric and completion are constructed. They are not the upload's unsaturated image-stability topology. \emph{C9, C10, C11.} \\[5pt]
ET48 & Universal claim rejected & The corrected constructive theory has closed written proofs. This is not a proof that the original three prose axioms force every later structure. \emph{C1--C14.} \\[5pt]
\end{longtable}
\normalsize

\end{document}
